% Provenance (never printed): see sources/outline_sources.md and sources/book1/BRIEFS.md (g8/03).
\chapter{Files and Data Tables}\label{ch:g8:files-and-data-tables}

The library of Brenmoor, the town of this series, lent books 400 times in
2025. Each loan is one line of a \omterm{def:g8:files-and-data-tables:text-file}{text file}: which book, which member, when it
went out, when it was due, when it came back. With a few lines of Python, that
\omterm{def:g3:files-and-folders:file}{file} answers questions that would take an afternoon by hand: which books are
still out, who borrows the most, which loans came back late.

\begin{recall}
A \omterm{def:g3:files-and-folders:file}{file} keeps data in \omterm{def:g3:inside-the-computer:storage}{storage}, under a name with an \omterm{def:g3:files-and-folders:extension}{extension}, in a \omterm{def:g3:files-and-folders:folder}{folder}
(\cref{ch:g3:files-and-folders}); \omterm{def:g3:inside-the-computer:storage}{storage} keeps it when the \omterm{def:g1:what-is-a-computer:computer}{computer} is off,
unlike the \omterm{def:g3:inside-the-computer:memory}{memory} (\cref{ch:g4:memory-and-storage}). A record is a \omterm{def:g8:dictionaries:dictionary}{dictionary}
whose keys are field names (\cref{ch:g8:dictionaries}). \code{split} cuts a
\omterm{def:g7:strings:string}{string} and \code{strip} removes the spaces at its ends
(\cref{ch:g7:strings}).
\end{recall}

\section{Text files}

\begin{definition}[Text file]\label{def:g8:files-and-data-tables:text-file}
A \emph{text file}\index{text file} is a \omterm{def:g3:files-and-folders:file}{file} that holds \omterm{def:g7:strings:string}{characters}, read and
written as text, line after line.
\end{definition}

\begin{definition}[End of line]\label{def:g8:files-and-data-tables:newline}
Each line of a \omterm{def:g8:files-and-data-tables:text-file}{text file} ends with an invisible \omterm{def:g7:strings:string}{character}, the \emph{end of
line}\index{end of line}, written \texttt{"\textbackslash n"} in Python.
\end{definition}

\code{open(name, "w", encoding="utf-8")} opens a \omterm{def:g3:files-and-folders:file}{file} for writing (it is
created, or emptied if it exists); \code{open(name, encoding="utf-8")} opens
it for reading. \code{encoding="utf-8"} says how the \omterm{def:g7:strings:string}{characters} are stored as
\omterm{def:g4:memory-and-storage:byte}{bytes}. The \code{with} block closes the \omterm{def:g3:files-and-folders:file}{file} at its end, so that everything
written is safely in \omterm{def:g3:inside-the-computer:storage}{storage}.

\omcode{code/grade-8/03-files-and-data-tables/listings/write_notes.py}{2}{12}{Writing
a \omterm{def:g3:files-and-folders:file}{file}, then reading it back twice.}\label{lst:g8:files-and-data-tables:notes}

\Needspace*{10\baselineskip}
\omoutput{code/grade-8/03-files-and-data-tables/outputs/write_notes.txt}{The
blank lines come from the \omterm{def:g8:files-and-data-tables:newline}{end of line} kept in each
line.}\label{out:g8:files-and-data-tables:notes}

The first \omterm{def:g4:algorithms-repeat-and-decide:loop}{loop} prints each line with its own \omterm{def:g8:files-and-data-tables:newline}{end of line}, and \code{print}
adds another one: hence the blank lines. \code{strip} removes it, and
\code{len} shows that it counts as a \omterm{def:g7:strings:string}{character}: \code{"Buy milk"} has 8
letters and spaces, the line has 9 \omterm{def:g7:strings:string}{characters}. After the \omterm{def:g4:programs-and-operating-system:program}{program} ends, the
\omterm{def:g3:files-and-folders:file}{file} \code{notes.txt} is still there, in \omterm{def:g3:inside-the-computer:storage}{storage}.

\section{A table saved as text}

\begin{definition}[Data table, CSV file, header row]\label{def:g8:files-and-data-tables:table}\label{def:g8:files-and-data-tables:csv}\label{def:g8:files-and-data-tables:header-row}
A \emph{data table}\index{data table} is a \omterm{def:g7:lists:list}{list} of records that all have the
same fields. A \emph{CSV file}\index{CSV file} (for \emph{comma-separated
\omterm{def:g6:variables-and-input:value}{values}}) keeps a data table as text: one record per line, the fields
separated by commas. Its first line, the \emph{header row}\index{header row},
gives the names of the fields.
\end{definition}

% ledger: g8:csv-rfc4180
The CSV format was written down in 2005, in a public document that the whole
internet uses as a reference. Here are the first lines of the library's book
table, \code{books.csv}:

\Needspace*{8\baselineskip}
\omcode{code/grade-8/03-files-and-data-tables/listings/read_lines.py}{1}{5}{The
first lines of a \omterm{def:g8:files-and-data-tables:table}{CSV file}.}\label{lst:g8:files-and-data-tables:read-lines}

\Needspace*{9\baselineskip}
\omoutput{code/grade-8/03-files-and-data-tables/outputs/read_lines.txt}{80 books,
plus the \omterm{def:g8:files-and-data-tables:table}{header row}.}\label{out:g8:files-and-data-tables:read-lines}

\begin{omfigure}
\begin{tikzpicture}[x=1cm, y=0.55cm, font=\scriptsize\ttfamily]
  \node[draw=black!50, fill=ocYellow!15, align=left, anchor=north west, inner sep=4pt] at (0,0) {book\_id,title,author,year,genre,pages\\ 1,The Silent Code,Fatou Dalewell,2002,mystery,65\\ 2,The Paper Forest,Zoe Stoneton,1996,sport,133};
  \draw[-{Stealth[length=2mm]}, thick] (3.3,-2.0) -- (3.3,-2.8);
  \foreach \h/\x/\w in {book\_id/0/1.1, title/1.1/2.0, author/3.1/1.9, year/5.0/0.7, genre/5.7/1.0, pages/6.7/0.8} {
    \draw[black!60, fill=ocInk!80] (\x,-4) rectangle ++(\w,1);
    \node[white] at (\x+\w/2,-3.5) {\h};
  }
  \foreach \pa/\pb/\pc/\pd/\pe/\pf [count=\i from 1] in {1/The Silent Code/Fatou Dalewell/2002/mystery/65, 2/The Paper Forest/Zoe Stoneton/1996/sport/133} {
    \foreach \v/\x/\w in {\pa/0/1.1, \pb/1.1/2.0, \pc/3.1/1.9, \pd/5.0/0.7, \pe/5.7/1.0, \pf/6.7/0.8} {
      \draw[black!60, fill=ocBlue!8] (\x,-4-\i) rectangle ++(\w,1);
      \node[font=\tiny\ttfamily] at (\x+\w/2,-3.5-\i) {\v};
    }
  }
\end{tikzpicture}

{\small A \omterm{def:g8:files-and-data-tables:table}{CSV file} as text (above) and the table it means (below): the \omterm{def:g8:files-and-data-tables:table}{header
row} names the fields, each other line is one record.}
\end{omfigure}

A line can be cut at its commas with \code{split(",")}. But a field may itself
hold a comma, such as the title \emph{Salt, Sand and Sea}; the CSV format then
puts the field between double quotes, and \code{split} cuts it in two. The
\code{csv} module of Python knows the rule:

\Needspace*{12\baselineskip}
\omcode{code/grade-8/03-files-and-data-tables/listings/split_line.py}{1}{9}{\code{split}
against the \code{csv} module.}\label{lst:g8:files-and-data-tables:split}

\Needspace*{7\baselineskip}
\omoutput{code/grade-8/03-files-and-data-tables/outputs/split_line.txt}{Seven
fields instead of six on the second line; the \code{csv} module gets six.}\label{out:g8:files-and-data-tables:split}

\section{Loading a table as a list of records}

\code{csv.DictReader} reads a \omterm{def:g8:files-and-data-tables:table}{CSV file} line by line and turns each line into a
\omterm{def:g8:dictionaries:dictionary}{dictionary}, whose keys come from the \omterm{def:g8:files-and-data-tables:table}{header row}: a record, as in
\cref{ch:g8:dictionaries}. \code{list(...)} collects them into a \omterm{def:g7:lists:list}{list}.

\Needspace*{13\baselineskip}
\omcode{code/grade-8/03-files-and-data-tables/listings/load_books.py}{1}{10}{The
book table as a \omterm{def:g7:lists:list}{list} of records.}\label{lst:g8:files-and-data-tables:load}

\Needspace*{8\baselineskip}
\omoutput{code/grade-8/03-files-and-data-tables/outputs/load_books.txt}{Every
\omterm{def:g6:variables-and-input:value}{value} read from a \omterm{def:g3:files-and-folders:file}{file} is text, a \code{str}.}\label{out:g8:files-and-data-tables:load}

A \omterm{def:g3:files-and-folders:file}{file} holds \omterm{def:g7:strings:string}{characters}, so \code{"65"} is two \omterm{def:g7:strings:string}{characters}, not a number:
\code{int} or \code{float} converts it before any calculation.

\begin{omfigure}
\begin{tikzpicture}[font=\scriptsize\ttfamily]
  \node[draw=black!50, fill=ocYellow!15, inner sep=4pt] (l) at (0,0) {1,The Silent Code,Fatou Dalewell,2002,mystery,65};
  \node[draw=black!50, fill=ocBlue!10, align=left, inner sep=4pt] (d) at (0,-2.1) {\{'book\_id': '1', 'title': 'The Silent Code',\\ \ 'author': 'Fatou Dalewell', 'year': '2002',\\ \ 'genre': 'mystery', 'pages': '65'\}};
  \draw[-{Stealth[length=2mm]}, thick] (l) -- node[right, font=\footnotesize\sffamily] {DictReader, with the header row's names as keys} (d);
\end{tikzpicture}

{\small One line of the \omterm{def:g3:files-and-folders:file}{file} becomes one \omterm{def:g8:dictionaries:dictionary}{dictionary}; every \omterm{def:g6:variables-and-input:value}{value} is still
text.}
\end{omfigure}

\section{Questions about a table}

Most questions about a table are a walk through its records (\cref{ch:g7:lists}):
selecting the records that pass a test, counting them, keeping the largest,
adding up a column. Which loans are still out? Those whose return date is
empty:

\omcode{code/grade-8/03-files-and-data-tables/listings/still_out.py}{1}{12}{Selecting
the loans still out.}\label{lst:g8:files-and-data-tables:still-out}

\Needspace*{8\baselineskip}
\omoutput{code/grade-8/03-files-and-data-tables/outputs/still_out.txt}{The first
three of them.}\label{out:g8:files-and-data-tables:still-out}

\begin{omfigure}
\begin{tikzpicture}[font=\footnotesize]
  \draw[thick, fill=ocBlue!10] (-2.4,0) -- (2.4,0) -- (0.7,-1.6) -- (0.7,-2.2) -- (-0.7,-2.2) -- (-0.7,-1.6) -- cycle;
  \node at (0,-0.5) {400 loans};
  \node[anchor=west, font=\scriptsize] at (2.6,-1.2) {test: \code{return_date == ""}};
  \node at (0,-1.9) {22};
  \node[anchor=west, align=left] at (2.6,-0.5) {every record is tested};
  \node[anchor=west, align=left] at (2.6,-1.9) {those that pass are kept};
\end{tikzpicture}

{\small Selecting rows: a test on each record, as a funnel.}
\end{omfigure}

The longest book keeps the best record so far, comparing pages as numbers:

\Needspace*{13\baselineskip}
\omcode{code/grade-8/03-files-and-data-tables/listings/longest_book.py}{1}{10}{The
longest book.}\label{lst:g8:files-and-data-tables:longest}

\Needspace*{5\baselineskip}
\omoutput{code/grade-8/03-files-and-data-tables/outputs/longest_book.txt}{513
pages.}\label{out:g8:files-and-data-tables:longest}

The weather table holds one row per day and per station. The mean temperature
at Hilltop adds the temperatures and divides by their number, as in a
mathematics lesson. But some \omterm{def:g6:variables-and-input:value}{values} are missing: the field is empty, and
\code{float("")} would stop the \omterm{def:g4:programs-and-operating-system:program}{program}, so the \omterm{def:g4:programs-and-operating-system:program}{program} skips them and counts
them.

\omcode{code/grade-8/03-files-and-data-tables/listings/mean_temp.py}{1}{17}{The
mean temperature at Hilltop, without the missing
\omterm{def:g6:variables-and-input:value}{values}.}\label{lst:g8:files-and-data-tables:mean}

\Needspace*{6\baselineskip}
\omoutput{code/grade-8/03-files-and-data-tables/outputs/mean_temp.txt}{365 days,
3 of them without a temperature.}\label{out:g8:files-and-data-tables:mean}

\begin{remark}[Data with holes]\label{rem:g8:files-and-data-tables:holes}
Real tables have mistakes: \omterm{def:g6:variables-and-input:value}{values} missing, \omterm{def:g6:variables-and-input:value}{values} typed twice, numbers that
mean ``the sensor failed'' rather than a measurement. The weather table has
all three, and the weekend problem finds them. This chapter only skips the
empty \omterm{def:g6:variables-and-input:value}{values}; cleaning a table properly is a topic of the next book.
\end{remark}

A \omterm{def:g4:programs-and-operating-system:program}{program} can also write its answers to a \omterm{def:g3:files-and-folders:file}{file}, as a report that stays after
it ends:

\omcode{code/grade-8/03-files-and-data-tables/listings/report.py}{2}{19}{A report
written to a \omterm{def:g3:files-and-folders:file}{file}, then printed.}\label{lst:g8:files-and-data-tables:report}

\Needspace*{8\baselineskip}
\omoutput{code/grade-8/03-files-and-data-tables/outputs/report.txt}{The text of
\code{report.txt}.}\label{out:g8:files-and-data-tables:report}

\begin{omfigure}
\begin{tikzpicture}
  \begin{axis}[width=11cm, height=4.6cm, ymin=0, xmin=0.5, xmax=12.5,
      xtick={1,...,12}, xticklabels={J,F,M,A,M,J,J,A,S,O,N,D},
      ylabel={$^\circ$C}, label style={font=\footnotesize}, tick label style={font=\footnotesize},
      axis x line*=bottom, axis y line*=left]
    \addplot[thick, ocRed, mark=*, mark size=1.5pt] table[x=month, y=temp] {code/grade-8/03-files-and-data-tables/figures/out/hilltop_monthly.dat};
  \end{axis}
\end{tikzpicture}

{\small The mean temperature of each month at Hilltop in 2025, computed from
the weather table by a script of the chapter.}
\end{omfigure}

\section{The same table in a spreadsheet}

\begin{definition}[Spreadsheet and formula]\label{def:g8:files-and-data-tables:spreadsheet}
A \emph{spreadsheet}\index{spreadsheet} is a \omterm{def:g4:programs-and-operating-system:program}{program} that shows a table as a
grid of cells, with lettered columns and numbered rows. A cell can hold a
\omterm{def:g6:variables-and-input:value}{value} or a \emph{formula}\index{formula (of a spreadsheet)}, a calculation on
other cells whose result the cell shows and keeps up to date.
\end{definition}

A \omterm{def:g8:files-and-data-tables:table}{CSV file} opens in any \omterm{def:g8:files-and-data-tables:spreadsheet}{spreadsheet} \omterm{def:g4:programs-and-operating-system:program}{program}: each field lands in a cell. A
\omterm{def:g8:files-and-data-tables:spreadsheet}{formula} such as \code{=SUM(F2:F81)} adds the pages of the 80 books; clicking a
column header sorts the rows by that column, and a filter hides the rows that
do not pass a test, like the selection of still-out loans.

\begin{omfigure}
\begin{tikzpicture}[x=1cm, y=0.5cm, font=\scriptsize]
  \foreach \c/\x/\w in {A/0.5/0.6, B/1.1/3.4, C/4.5/1.0, D/5.5/1.0} {
    \draw[black!50, fill=black!10] (\x,0) rectangle ++(\w,1); \node at (\x+\w/2,0.5) {\c};
  }
  \foreach \r in {1,...,5} {
    \draw[black!50, fill=black!10] (0,-\r) rectangle ++(0.5,1); \node at (0.25,-\r+0.5) {\r};
  }
  \foreach \v/\x/\w [count=\r from 1] in {title/1.1/3.4, The Paper River/1.1/3.4, A Year by the Harbour/1.1/3.4, Notes on the Comet/1.1/3.4} {
    \draw[black!40, fill=white] (\x,-\r) rectangle ++(\w,1); \node[anchor=west] at (\x+0.05,-\r+0.5) {\v};
  }
  \foreach \v [count=\r from 1] in {pages, 513, 506, 502} {
    \draw[black!40, fill=white] (4.5,-\r) rectangle ++(1,1); \node[anchor=east] at (5.45,-\r+0.5) {\v};
  }
  \foreach \r in {1,...,5} {
    \draw[black!40, fill=white] (0.5,-\r) rectangle ++(0.6,1);
    \draw[black!40, fill=white] (5.5,-\r) rectangle ++(1,1);
  }
  \draw[black!40, fill=white] (1.1,-5) rectangle ++(3.4,1); \node[anchor=west] at (1.15,-4.5) {total};
  \draw[ocBlue, thick, fill=ocBlue!10] (4.5,-5) rectangle ++(1,1); \node[anchor=east] at (5.45,-4.5) {1521};
  \node[anchor=west, font=\scriptsize\ttfamily, ocBlue] (f) at (7.0,-4.5) {=SUM(C2:C4)};
  \draw[-{Stealth[length=1.5mm]}, ocBlue] (f) -- (5.55,-4.5);
  \node[anchor=west, ocRed] at (7.0,-2) {sorted by pages, largest first};
  \draw[-{Stealth[length=1.5mm]}, ocRed] (7.0,-2) -- (5.55,-2);
\end{tikzpicture}

{\small A \omterm{def:g8:files-and-data-tables:spreadsheet}{spreadsheet} grid with the three longest books: lettered columns,
numbered rows, a column sorted from the largest, and a \omterm{def:g8:files-and-data-tables:spreadsheet}{formula} cell whose
\omterm{def:g6:variables-and-input:value}{value} is the sum of three cells.}
\end{omfigure}

A \omterm{def:g8:files-and-data-tables:spreadsheet}{spreadsheet} keeps every step visible and is quick for one table. A \omterm{def:g4:programs-and-operating-system:program}{program}
does the same steps the same way every time, and on a thousand \omterm{def:g3:files-and-folders:file}{files} as
easily as on one. Both read the same \omterm{def:g8:files-and-data-tables:table}{CSV file}.

\section{Exercises}

\begin{exercise}[$\star$]\label{exo:g8:files-and-data-tables:1}
\code{books.csv} is replaced by a \omterm{def:g3:files-and-folders:file}{file} of three lines. What does
\cref{lst:g8:files-and-data-tables:read-lines} print?
\end{exercise}

\begin{exercise}[$\star$]\label{exo:g8:files-and-data-tables:2}
What does \code{"a,b,,c".split(",")} give? How many fields is that?
\end{exercise}

\begin{exercise}[$\star$]\label{exo:g8:files-and-data-tables:3}
What is the type of \code{book["year"]} for a record read with
\code{DictReader}? How do you get the year as a number?
\end{exercise}

\begin{exercise}[$\star$]\label{exo:g8:files-and-data-tables:4}
Write a \omterm{def:g7:functions:function}{function} \code{count_genres(books)} that returns a \omterm{def:g8:dictionaries:dictionary}{dictionary} giving
the number of books of each genre.
\end{exercise}

\begin{exercise}[$\star\star$]\label{exo:g8:files-and-data-tables:5}
This \omterm{def:g7:functions:function}{function} stops with a \code{TypeError}. Why? Fix it.

\omcode{code/grade-8/03-files-and-data-tables/exercises/pages_bug.py}{1}{5}{The
total of a column.}\label{lst:g8:files-and-data-tables:pages-bug}
\end{exercise}

\begin{exercise}[$\star\star$]\label{exo:g8:files-and-data-tables:6}
% ledger: g8:python-open-encoding
The town's table of places names a \emph{Lighthouse Café}. Opened with
\code{open("places.csv")}, without \code{encoding}, the \omterm{def:g3:files-and-folders:file}{file} reads well on one
\omterm{def:g1:what-is-a-computer:computer}{computer} but shows \texttt{Caf\~{A}\textcopyright} on another. Why? How do you make the \omterm{def:g4:programs-and-operating-system:program}{program}
read it the same way everywhere?
\end{exercise}

\begin{exercise}[$\star\star$]\label{exo:g8:files-and-data-tables:7}
Write a \omterm{def:g7:functions:function}{function} \code{most_loans(loans)} that returns the member who borrowed
the most books, and how many.
\end{exercise}

\begin{exercise}[$\star\star$]\label{exo:g8:files-and-data-tables:8}
Write a \omterm{def:g7:functions:function}{function} \code{late_returns(loans)} that counts the loans returned
after their due date. The dates are written \code{YYYY-MM-DD}, as text: why
can they be compared with \code{>} without converting them?
\end{exercise}

\begin{exercise}[$\star\star$]\label{exo:g8:files-and-data-tables:9}
How many lines does a \omterm{def:g8:files-and-data-tables:table}{CSV file} of 400 records have? And how many end-of-line
\omterm{def:g7:strings:string}{characters}, if every line ends with one?
\end{exercise}

\begin{exercise}[$\star\star\star$]\label{exo:g8:files-and-data-tables:10}
Write a \omterm{def:g7:functions:function}{function} \code{write_genre(books, genre, filename)} that writes the
books of one genre to a new \omterm{def:g8:files-and-data-tables:table}{CSV file} with the same \omterm{def:g8:files-and-data-tables:table}{header row}. Use
\code{csv.DictWriter}, the writing twin of \code{DictReader}.
\end{exercise}

\begin{exercise}[$\star\star\star$]\label{exo:g8:files-and-data-tables:11}
Comparing dates as text works for \code{2025-01-09} and \code{2025-01-10}.
Show with an example that it fails for dates written \code{DD-MM-YYYY}, and
explain why the order \code{YYYY-MM-DD} makes it work.
\end{exercise}

\begin{exercise}[$\star\star\star$]\label{exo:g8:files-and-data-tables:12}
Write a \omterm{def:g7:functions:function}{function} \code{rainiest_week(rows, station)} that finds the 7 days in a
row with the most rain at a station, and returns the first day and the total.
Count a date only once if it appears twice.
\end{exercise}

\section{Problem: A Year of Weather}

\begin{problem}[{Weekend problem --- a year of weather in Brenmoor}]\label{pb:g8:files-and-data-tables:1}
The \omterm{def:g3:files-and-folders:file}{file} \code{weather.csv} holds the readings of 2025 at the town's two
weather stations, Hilltop and Riverside: the date, the station, the
temperature in degrees, the rain in millimetres and the wind in kilometres
per hour.

\medskip\noindent\textbf{Part I --- By hand.}
\begin{enumerate}
  \item The first two lines after the \omterm{def:g8:files-and-data-tables:table}{header row} are
    \code{2025-01-01,Hilltop,3.9,0.0,15.3} and
    \code{2025-01-01,Riverside,4.3,0.0,10.4}. What do they say?
  \item How many lines should the \omterm{def:g3:files-and-folders:file}{file} have for two stations over a year of
    365 days, \omterm{def:g8:files-and-data-tables:table}{header row} included?
  \item The \omterm{def:g3:files-and-folders:file}{file} has 732 lines. What might explain the difference?
\end{enumerate}

\medskip\noindent\textbf{Part II --- Hilltop.}

\omcode{code/grade-8/03-files-and-data-tables/listings/weather_year.py}{1}{21}{Rain
per month and warm days at Hilltop.}\label{lst:g8:files-and-data-tables:weather-year}

\begin{enumerate}[resume]
  \item What are the keys of the \omterm{def:g8:dictionaries:dictionary}{dictionary} \code{rain}, and how are they cut
    out of the date?
  \item Why does the test of a warm day check \code{row["temp_c"] != ""}
    first?
  \item Run the \omterm{def:g4:programs-and-operating-system:program}{program}. Which month is the wettest, and how many days are
    above 20 degrees?
  \item Why does the search for the wettest month start from \code{"01"}?
\end{enumerate}

\medskip\noindent\textbf{Part III --- Riverside's holes.}

\omcode{code/grade-8/03-files-and-data-tables/listings/defects.py}{1}{30}{Looking
for the defects of Riverside's readings.}\label{lst:g8:files-and-data-tables:defects}

\begin{enumerate}[resume]
  \item Which three kinds of defect does the \omterm{def:g4:programs-and-operating-system:program}{program} look for?
  \item Run it. How many of each are there?
  \item A thermometer at Riverside never shows $-99$ degrees. What could that
    \omterm{def:g6:variables-and-input:value}{value} mean?
  \item By how much do the two codes pull the year's mean temperature down?
  \item Explain why two \omterm{def:g6:variables-and-input:value}{values} out of more than 360 make such a difference.
  \item State the result.
\end{enumerate}
\end{problem}
